% Generated draft: requires mathematical and source review.
\documentclass[11pt]{article}
\usepackage[margin=25mm]{geometry}
\usepackage{fontspec}
\setmainfont{Latin Modern Roman}
\usepackage{amsmath,amssymb,mathtools,unicode-math}
\setmathfont{Latin Modern Math}
\usepackage{enumitem,needspace,xurl,hyperref}
\setlength{\parindent}{0pt}
\setlength{\parskip}{5pt}
\setlength{\emergencystretch}{3em}
\newcommand{\N}{\mathbb{N}}
\newcommand{\Z}{\mathbb{Z}}
\newcommand{\Q}{\mathbb{Q}}
\newcommand{\R}{\mathbb{R}}
\newcommand{\C}{\mathbb{C}}
\newcommand{\F}{\mathbb{F}}
\newcommand{\A}{\mathbb{A}}
\newcommand{\PP}{\mathbb{P}}
\newcommand{\E}{\mathbb{E}}
\newcommand{\eps}{\varepsilon}
\newcommand{\dd}{\,\mathrm{d}}
\begin{document}
\section*{UnsolvedMath ID: 30004614}
\textbf{Source problem code:} catalogue entry
\subsection{Definitions and mathematical statement}
The supplied record asks two questions about vertex algebras and their partition functions. It does not specify the class of vertex algebras, the meaning of arising from a quadratic form, or the spaces and normalizations of partition functions. Consequently, the following is a quantified formulation scheme whose unspecified data require mathematical review.

Let \(\mathcal V\) denote the intended class of vertex algebras, and let \(\mathcal Q\subseteq\mathcal V\) denote those described in the source as arising from quadratic forms. Neither class is defined in the supplied record. For \(V\in\mathcal V\) and each relevant genus \(g\), write \(Z_g(V)\) for its partition function. Write \(\mathcal S_g(V)\) for the space of Siegel modular forms with the data needed for comparison with \(Z_g(V)\), and write
\[
\rho_{g,V}^{*}\colon\mathcal S_g(V)\longrightarrow\mathcal T_g(V)
\]
for the intended pullback map into a space containing \(Z_g(V)\). Here these symbols are placeholders: the moduli spaces, map, weights, levels, representations, line bundles, and boundary conditions are not specified. For an ordinary function \(F\), pullback along a map \(\rho\) means composition \(\rho^{*}F=F\circ\rho\); a formulation using sections also requires compatible line-bundle identifications.

\textbf{Question 1: eventual failure to be a pullback.} Once these data have been fixed, is it true that
\[
\forall V\in\mathcal V\setminus\mathcal Q\quad
\exists G(V)\in\Z_{\geq 0}\quad
\forall g\in\Z_{\geq 0},\ g\geq G(V):\qquad
Z_g(V)\notin\operatorname{im}(\rho_{g,V}^{*})?
\]
This interprets sufficiently large genus as every genus beyond a threshold that may depend on \(V\), assuming the partition functions and pullback maps exist in those genera. The source does not state a uniform threshold or specify existence hypotheses.

\textbf{Question 2: determination by partition functions.} Let \(I\) be the intended set of genera, and suppose compatible comparison data make equality of the indexed partition functions meaningful. Is the assignment
\[
[V]\longmapsto\bigl(Z_g(V)\bigr)_{g\in I}
\]
injective on isomorphism classes in the intended domain? Equivalently, does
\[
\forall V,W\in\mathcal V,\qquad
\left(\forall g\in I,\ Z_g(V)=Z_g(W)\right)
\Longrightarrow V\cong W
\]
hold? Here \(V\cong W\) means an isomorphism preserving the vertex-algebra structure and any additional structure required in the intended category. The record does not specify that additional structure, the equality convention for partition functions, or whether this second question concerns all of \(\mathcal V\) or only \(\mathcal V\setminus\mathcal Q\). Thus neither question is supplied with enough hypotheses for an unconditional precise conjecture.
\subsection{Short English statement}
The first question asks whether every vertex algebra outside the unspecified quadratic-form class eventually has partition functions that cannot be obtained by pulling back Siegel modular forms. The second asks whether the full indexed family of partition functions determines the vertex algebra up to isomorphism. Both questions require additional hypotheses and comparison conventions. The supplied partially solved label is dated curation, not an independently verified mathematical conclusion.
\subsection{Sources}
\begin{itemize}
\item Local UnsolvedMath registry: permanent ID \(30004614\), source problem code \texttt{catalogue entry}, title \emph{Vertex Algebra Partition Functions and Siegel Modular Forms}, in the Oberwolfach Reports Open Problems collection. The original, clean, and statement fields give the same two questions. The supplied proposed year is 2021; no individual proposer is named.
\item Supplied source citation: \emph{Moduli spaces and Modular forms (2022)}, Oberwolfach Reports, pp.\ 301--359, DOI \texttt{10.4171/owr/2021/6}. \url{https://doi.org/10.4171/owr/2021/6}. The background describes Oberwolfach Reports as documenting mathematical workshops and their open-problem discussions. No specific page or passage for these questions is supplied.
\item Supplied literature reference: Carpi--Codogni, \emph{Vertex operator algebras, partition functions and Teichmuller modular forms}. \url{https://arxiv.org/abs/2605.26972}.
\item Supplied literature reference: \emph{Oberwolfach Report 6/2021}. \url{https://doi.org/10.4171/owr/2021/6}.
\item Literature assessment dated exactly \texttt{2026-08-22}: the supplied curation labels the problem partially solved and states that recent work constructs Teichmuller partition functions and proves non-pullback results in broad cases, while the full non-pullback conjecture and reconstruction from all partition functions remain open. This is supplied dated curation derived from the OWR record, not an independently verified assessment; the source explicitly requests independent verification of time-sensitive status and cited claims.
\item The registry labels the statement \texttt{corrected\_verified} and reports a manual semantic, typographic, and LaTeX review asserting that the clean statement is self-contained and preserves its target. These are supplied editorial assertions; this draft does not independently verify them or establish mathematical correctness.
\end{itemize}
\end{document}
